\begin{abstract}
Autonomous research systems are usually evaluated on whether they produce an
answer, execute a program, or reproduce a published result. These evaluations
do not establish whether a system can turn its own failed experiment into a
new, falsifiable mechanism and preserve the resulting evidence without
post-result adaptation. We study that narrower capability under a
preregistered, fully local protocol. \method{} binds each hypothesis to the
hash of a parent negative result, freezes code and adjudication before unseen
evaluation, and opens a new round rather than tuning after reveal. We first
exercise this protocol in two real equation-discovery rounds. Both candidates
passed development screening, both failed their frozen unseen confidence
gate, and both were retained as negative results. The resulting pivot then
motivates a systems study over \TaskCount{} tasks from two task families,
\SeedCount{} seeds, three controllers, and four component ablations
(\CellCount{} total cells). The full loop attains task success
\FullLoopSuccess{} versus \ExecuteOnceSuccess{} for execute-once, a paired
gain of \PairedGain{} with a \BootstrapResamples{}-resample bootstrap
confidence interval of [\PairedCILower{}, \PairedCIUpper{}]. It recovers
\FullLoopRecovery{} of initially negative cases, exactly reproduces
\FullLoopReproduction{} of evaluated cells, and emits \FullLoopUnsupported{}
unsupported scientific claims under the frozen evaluator. These results
support a systems claim about evidence-bound recovery in this controlled
task matrix, not a claim of general autonomous discovery or a new
equation-discovery method. We release the complete hash-linked campaign,
negative results, source evidence, paper source, and clean-directory
reproduction entry point for human review.
\end{abstract}
